\documentclass{chsmc}
\chsmcsetup{
	year = 2011,
	part = 1,
	type = B,
	show-answers = false,
}
\begin{document}

\section{填空题：本大题共 8 小题，每小题 8 分，共 64 分.}

\begin{enumerate}
	\item 设等差数列 $\{a_n\}$ 的前 $n$ 项和为 $S_n$，若 $S_{2010} - S_1 = 1$，
		则 $S_{2011} =$ \blankline
	\item 已知复数 $z$ 的模为 $1$，若 $z = z_1$ 和 $z = z_2$ 时
		$|z + 1 + \ii|$ 分别取得最大值和最小值，则 $z_1 - z_2 =$ \blankline
	\item 设 $a$, $b$ 为正实数，$\dfrac{1}{a} + \dfrac{1}{b} \leqslant 2\sqrt{2}$，
		$(a-b)^2 = 4(ab)^3$，则 $\log_a b =$ \blankline
	\item 把扑克牌中的 A, 2, $\ldots$, J, Q, K 分别看作数字 $1$, $2$, $\ldots$, $11$, $12$, $13$.
		现将一幅扑克牌中的黑桃、红桃各 $13$ 张放在一起，从中随机取出 $2$ 张牌，
		其花色相同且两个数的积是完全平方数的概率为 \blankline
	\item 已知正方体 $ABCD$-$A_1B_1C_1D_1$ 的棱长为 $6$，$M$, $N$ 分别是 $BB_1$, $B_1C_1$
		上的点，$B_1M = B_1N = 2$，$S$, $P$ 分别是线段 $AD$, $MN$ 的中点，
		则异面直线 $SP$ 与 $AC_1$ 的距离为 \blankline
	\item 在 $\triangle ABC$ 中，$E$, $F$ 分别是 $AC$, $AB$ 的中点，
		$AB = \dfrac{2}{3}AC$. 若 $\dfrac{BE}{CF} < t$ 恒成立，则 $t$ 的最小值为 \blankline
	\item 抛物线 $y^2 = 2p\left(x - \dfrac{p}{2}\right)$（$p > 0$）上动点 $A$ 到点 $B(3,0)$
		的距离的最小值记为 $d(p)$，满足 $d(p) = 2$ 的所有实数 $p$ 的和为 \blankline
\end{enumerate}

\section{解答题：本大题共 3 个小题，满分 56 分. 解答应写出文字说明、证明过程或演算步骤.}

\begin{enumerate}[resume]
	\item (本题满分 16 分) 已知实数 $x$, $y$, $z$ 满足：$x \geqslant y \geqslant z$，
		$x + y + z = 1$，$x^2 + y^2 + z^2 = 3$. 求实数 $x$ 的取值范围.
	\item (本题满分 20 分) 设数列 $\{a_n\}$ 的前 $n$ 项和 $S_n$ 组成的数列满足
		$S_n + S_{n+1} + S_{n+2} = 6n^2 + 9n + 7$（$n \geqslant 1$）. 已知 $a_1 = 1$，
		$a_2 = 5$，求数列 $\{a_n\}$ 的通项公式.
	\item (本题满分 20 分) 已知 $A_1(x_1,y_1)$, $A_2(x_2,y_2)$, $A_3(x_3,y_3)$ 是抛物线
		$y^2 = 2px$（$p > 0$）上不同的三点，$\triangle A_1A_2A_3$ 有两边所在的直线与抛物线
		$x^2 = 2qy$（$q > 0$）相切，证明：对不同的 $i$, $j \in \{1,2,3\}$，
		$y_i y_j (y_i + y_j)$ 为定值.
\end{enumerate}

\end{document}
